\documentclass[11pt, a4paper]{article}\usepackage[]{graphicx}\usepackage[]{xcolor}
% maxwidth is the original width if it is less than linewidth
% otherwise use linewidth (to make sure the graphics do not exceed the margin)
\makeatletter
\def\maxwidth{ %
  \ifdim\Gin@nat@width>\linewidth
    \linewidth
  \else
    \Gin@nat@width
  \fi
}
\makeatother

\definecolor{fgcolor}{rgb}{0.345, 0.345, 0.345}
\newcommand{\hlnum}[1]{\textcolor[rgb]{0.686,0.059,0.569}{#1}}%
\newcommand{\hlsng}[1]{\textcolor[rgb]{0.192,0.494,0.8}{#1}}%
\newcommand{\hlcom}[1]{\textcolor[rgb]{0.678,0.584,0.686}{\textit{#1}}}%
\newcommand{\hlopt}[1]{\textcolor[rgb]{0,0,0}{#1}}%
\newcommand{\hldef}[1]{\textcolor[rgb]{0.345,0.345,0.345}{#1}}%
\newcommand{\hlkwa}[1]{\textcolor[rgb]{0.161,0.373,0.58}{\textbf{#1}}}%
\newcommand{\hlkwb}[1]{\textcolor[rgb]{0.69,0.353,0.396}{#1}}%
\newcommand{\hlkwc}[1]{\textcolor[rgb]{0.333,0.667,0.333}{#1}}%
\newcommand{\hlkwd}[1]{\textcolor[rgb]{0.737,0.353,0.396}{\textbf{#1}}}%
\let\hlipl\hlkwb

\usepackage{framed}
\makeatletter
\newenvironment{kframe}{%
 \def\at@end@of@kframe{}%
 \ifinner\ifhmode%
  \def\at@end@of@kframe{\end{minipage}}%
  \begin{minipage}{\columnwidth}%
 \fi\fi%
 \def\FrameCommand##1{\hskip\@totalleftmargin \hskip-\fboxsep
 \colorbox{shadecolor}{##1}\hskip-\fboxsep
     % There is no \\@totalrightmargin, so:
     \hskip-\linewidth \hskip-\@totalleftmargin \hskip\columnwidth}%
 \MakeFramed {\advance\hsize-\width
   \@totalleftmargin\z@ \linewidth\hsize
   \@setminipage}}%
 {\par\unskip\endMakeFramed%
 \at@end@of@kframe}
\makeatother

\definecolor{shadecolor}{rgb}{.97, .97, .97}
\definecolor{messagecolor}{rgb}{0, 0, 0}
\definecolor{warningcolor}{rgb}{1, 0, 1}
\definecolor{errorcolor}{rgb}{1, 0, 0}
\newenvironment{knitrout}{}{} % an empty environment to be redefined in TeX

\usepackage{alltt}

\usepackage[top = 0.6 in, bottom = 0.6 in, left = 0.8 in, right = 0.8 in]{geometry}
\usepackage{amsmath, amssymb, amsfonts}

\allowdisplaybreaks[4]

\usepackage{enumerate}
\usepackage{array}
\usepackage{multirow}
\usepackage{fontawesome5}
\usepackage{tasks}
\usepackage{bbding}
\usepackage{undertilde}
\usepackage{twemojis}
% how to use bull's eye ----- \scalebox{2.0}{\twemoji{bullseye}}
\usepackage{fontspec}
\usepackage{fancyhdr}
\usepackage{lastpage}
\usepackage{mathrsfs}
\usepackage{pifont}
\usepackage{hyperref}

\pagestyle{fancy}
\fancyhf{}
\fancyfoot[C]{Page \thepage\ of \pageref{LastPage}}
\renewcommand{\headrulewidth}{0pt}

\title{MSMS 408 : Practical 16}
\author{{\fontsize{25}{20} \benyorithafont Ananda Biswas} \\[1em] Exam Roll No. : 24419STC053}
\date{\today}

\newfontface\benyorithafont{Benyoritha-G334D.otf}

\newfontface\myfont{Myfont1-Regular.ttf}[LetterSpace=0.05em]

\newfontface\cbfont{CaveatBrush-Regular.ttf}
% how to use --- \myfont --write text here--

\newfontface\gloriahallelujahfont{GLORIAHALLELUJAH-REGULAR.TTF}
\IfFileExists{upquote.sty}{\usepackage{upquote}}{}
\begin{document}

\maketitle


\section*{\faArrowAltCircleRight[regular] \textcolor{blue}{Question}}

Consider a random variable $X \sim \mathscr{B}e(2, 5)$ with mean $E[X] = \dfrac{2}{7}$.

\begin{enumerate}[(i)]
\item Estimate the mean using Monte Carlo simulation and the antithetic variable technique. Compare the variances of the two techniques.

\item Estimate $E[e^{-X}]$ using Monte Carlo simulation and the antithetic variable technique. Compare the variances of the two techniques.

\item Estimate $E[X^2]$ using Monte Carlo simulation and the antithetic variable technique. Compare the variances of the two techniques.
\end{enumerate}


\section*{\faArrowAltCircleRight[regular] \textcolor{blue}{Theory}}

\begin{itemize}

\item To generate random numbers from $\mathscr{B}e_1$ distribution, we recall

$$X \sim U(0, 1) \Rightarrow - \ln X \sim \; \text{Exp}(1).$$

Also, $$\sum \limits_{i = 1}^{n} - \ln X_i \sim \; \text{Gamma}(\text{shape} = n, \text{scale} = 1) \qquad \text{and}$$

$$\dfrac{S_1}{S_1 + S_2} \sim \beta_1(m, n)$$
where $S_1$ and $S_2$ are independent Gamma distributions with shape parameters $m$ and $n$ respectively.


\item If $x_1, x_2, \ldots, x_n$ is a sample of size $n$ from $X$, Monte Carlo estimate of $\theta = E[X]$ is $$\widehat{\theta} = \dfrac{1}{n} \sum \limits_{i = 1}^{n} x_i.$$

Estimate of $\theta$ by using antithetic variable will be

$$\widehat{\theta} = \dfrac{1}{n} \sum \limits_{i = 1}^{n/2} [X_i + X_i']$$

where $X_i$ and $X_i'$ are negatively correlated random variables.


\item If $x_1, x_2, \ldots, x_n$ is a sample of size $n$ from $X$, Monte Carlo estimate of $\theta = E[e^{-X}]$ is $$\widehat{\theta} = \dfrac{1}{n} \sum \limits_{i = 1}^{n} e^{-x_i}.$$

Estimate of $\theta$ by using antithetic variable will be

$$\widehat{\theta} = \dfrac{1}{n} \sum \limits_{i = 1}^{n/2} [e^{-X_i} + e^{-X_i'}]$$

where $X_i$ and $X_i'$ are negatively correlated random variables; $g(x) = e^{-x}$ is a monotone decreasing function of $x$, consequently $e^{-X_i}$ and $e^{-X_i'}$ are negatively correlated random variables.


\item If $x_1, x_2, \ldots, x_n$ is a sample of size $n$ from $X$, Monte Carlo estimate of $\theta = E[X^2]$ is $$\widehat{\theta} = \dfrac{1}{n} \sum \limits_{i = 1}^{n} x_i^2.$$

Estimate of $\theta$ by using antithetic variable will be

$$\widehat{\theta} = \dfrac{1}{n} \sum \limits_{i = 1}^{n/2} [X_i^2 + X_i'^2]$$

where $X_i$ and $X_i'$ are negatively correlated random variables; $g(x) = x^2$ is a monotone increasing function of $x$, consequently $X_i^2$ and $X_i'^2$ are negativly correlated random variables.

\item Here

$$X_i = h(\utilde{u}), \qquad \& \qquad X_i' = h(\mathbf{1} - \utilde{u})$$

with

$$\utilde{u} = (u_1, u_2, \ldots, u_7), \qquad \text{and}$$

$$h(u_1, u_2, \ldots, u_7) = \dfrac{\sum \limits_{i = 1}^{2} - \ln u_i}{\sum \limits_{i = 1}^{2} - \ln u_i + \sum \limits_{i = 3}^{7} - \ln u_i}.$$

\end{itemize}

\section*{\faArrowAltCircleRight[regular] \textcolor{blue}{R Program}}






\begin{knitrout}
\definecolor{shadecolor}{rgb}{0.969, 0.969, 0.969}\color{fgcolor}\begin{kframe}
\begin{alltt}
\hldef{h} \hlkwb{<-} \hlkwa{function}\hldef{(}\hlkwc{u}\hldef{,} \hlkwc{m}\hldef{,} \hlkwc{n}\hldef{)\{}
  \hldef{u1} \hlkwb{<-} \hldef{u[}\hlnum{1}\hlopt{:}\hldef{m]}
  \hldef{u2} \hlkwb{<-} \hldef{u[(m} \hlopt{+} \hlnum{1}\hldef{)} \hlopt{:} \hldef{(m} \hlopt{+} \hldef{n)]}

  \hldef{x1} \hlkwb{<-} \hlopt{-}\hlkwd{sum}\hldef{(}\hlkwd{log}\hldef{(u1))}
  \hldef{x2} \hlkwb{<-} \hlopt{-}\hlkwd{sum}\hldef{(}\hlkwd{log}\hldef{(u2))}

  \hldef{x} \hlkwb{<-} \hldef{x1} \hlopt{/} \hldef{(x1} \hlopt{+} \hldef{x2)}
\hldef{\}}
\end{alltt}
\end{kframe}
\end{knitrout}

\begin{knitrout}
\definecolor{shadecolor}{rgb}{0.969, 0.969, 0.969}\color{fgcolor}\begin{kframe}
\begin{alltt}
\hldef{my.rbeta} \hlkwb{<-} \hlkwa{function}\hldef{(}\hlkwc{n}\hldef{,} \hlkwc{shape1}\hldef{,} \hlkwc{shape2}\hldef{,} \hlkwc{uniform.sample}\hldef{)\{}

  \hldef{s} \hlkwb{<-} \hlkwd{c}\hldef{()}

  \hlkwa{for} \hldef{(i} \hlkwa{in} \hlnum{1}\hlopt{:}\hldef{n) \{}
    \hldef{s[i]} \hlkwb{<-} \hlkwd{h}\hldef{(uniform.sample[i,], shape1, shape2)}
  \hldef{\}}

  \hlkwd{return}\hldef{(s)}
\hldef{\}}
\end{alltt}
\end{kframe}
\end{knitrout}

\begin{knitrout}
\definecolor{shadecolor}{rgb}{0.969, 0.969, 0.969}\color{fgcolor}\begin{kframe}
\begin{alltt}
\hldef{sample.sizes} \hlkwb{<-} \hlkwd{c}\hldef{(}\hlnum{20}\hldef{,} \hlnum{50}\hldef{,} \hlnum{80}\hldef{,} \hlnum{100}\hldef{)}
\end{alltt}
\end{kframe}
\end{knitrout}

\begin{enumerate}

\item[\scalebox{2.0}{\twemoji{keycap: 1}}]

\begin{knitrout}
\definecolor{shadecolor}{rgb}{0.969, 0.969, 0.969}\color{fgcolor}\begin{kframe}
\begin{alltt}
\hldef{theta.hat.monte.carlo} \hlkwb{<-} \hlkwd{c}\hldef{()}

\hlkwa{for} \hldef{(i} \hlkwa{in} \hlnum{1}\hlopt{:}\hlkwd{length}\hldef{(sample.sizes)) \{}

  \hldef{size} \hlkwb{<-} \hldef{sample.sizes[i]}

  \hldef{u} \hlkwb{<-} \hlkwd{matrix}\hldef{(}\hlkwd{runif}\hldef{(size} \hlopt{*} \hldef{(}\hlnum{2} \hlopt{+} \hlnum{5}\hldef{)),} \hlkwc{nrow} \hldef{= size,} \hlkwc{ncol} \hldef{= (}\hlnum{2} \hlopt{+} \hlnum{5}\hldef{))}

  \hldef{theta.hat.monte.carlo[i]} \hlkwb{<-} \hlkwd{mean}\hldef{(}\hlkwd{my.rbeta}\hldef{(size,} \hlnum{2}\hldef{,} \hlnum{5}\hldef{, u))}
\hldef{\}}
\end{alltt}
\end{kframe}
\end{knitrout}



\begin{knitrout}
\definecolor{shadecolor}{rgb}{0.969, 0.969, 0.969}\color{fgcolor}\begin{kframe}
\begin{alltt}
\hldef{theta.hat.antithetic} \hlkwb{<-} \hlkwd{c}\hldef{()}

\hlkwa{for} \hldef{(i} \hlkwa{in} \hlnum{1}\hlopt{:}\hlkwd{length}\hldef{(sample.sizes)) \{}
  \hldef{size} \hlkwb{<-} \hldef{sample.sizes[i]}

  \hldef{u} \hlkwb{<-} \hlkwd{matrix}\hldef{(}\hlkwd{runif}\hldef{((size} \hlopt{/} \hlnum{2}\hldef{)} \hlopt{*} \hldef{(}\hlnum{2} \hlopt{+} \hlnum{5}\hldef{)),} \hlkwc{nrow} \hldef{= size} \hlopt{/} \hlnum{2}\hldef{,} \hlkwc{ncol} \hldef{= (}\hlnum{2} \hlopt{+} \hlnum{5}\hldef{))}

  \hldef{x} \hlkwb{<-} \hlkwd{my.rbeta}\hldef{(size} \hlopt{/} \hlnum{2}\hldef{,} \hlnum{2}\hldef{,} \hlnum{5}\hldef{, u)}
  \hldef{x_dash} \hlkwb{<-} \hlkwd{my.rbeta}\hldef{(size} \hlopt{/} \hlnum{2}\hldef{,} \hlnum{2}\hldef{,} \hlnum{5}\hldef{,} \hlnum{1}\hlopt{-}\hldef{u)}

  \hldef{theta.hat.antithetic[i]} \hlkwb{<-} \hlkwd{mean}\hldef{(}\hlkwd{c}\hldef{(x, x_dash))}
\hldef{\}}
\end{alltt}
\end{kframe}
\end{knitrout}

\begin{knitrout}
\definecolor{shadecolor}{rgb}{0.969, 0.969, 0.969}\color{fgcolor}\begin{kframe}
\begin{alltt}
\hldef{df1} \hlkwb{<-} \hlkwd{data.frame}\hldef{(}\hlkwc{sample.size} \hldef{= sample.sizes,}
                  \hlkwc{theta.hat.mc} \hldef{= theta.hat.monte.carlo,}
                  \hlkwc{theta.hat.at} \hldef{= theta.hat.antithetic)}
\end{alltt}
\end{kframe}
\end{knitrout}


% Table created by stargazer v.5.2.3 by Marek Hlavac, Social Policy Institute. E-mail: marek.hlavac at gmail.com
% Date and time: Wed, May 20, 2026 - 02:27:17
\begin{table}[!htbp] \centering 
  \caption{Estimates from Two Different Techniques} 
  \label{Table 1} 
\begin{tabular}{@{\extracolsep{5pt}} ccc} 
\\[-1.8ex]\hline 
\hline \\[-1.8ex] 
sample.size & theta.hat.mc & theta.hat.at \\ 
\hline \\[-1.8ex] 
$20$ & $0.199437$ & $0.289805$ \\ 
$50$ & $0.287849$ & $0.286490$ \\ 
$80$ & $0.291264$ & $0.290699$ \\ 
$100$ & $0.283530$ & $0.293247$ \\ 
\hline \\[-1.8ex] 
\end{tabular} 
\end{table} 


True value is $\theta = E[X] = \dfrac{2}{7} = 0.2857143$.

\begin{knitrout}
\definecolor{shadecolor}{rgb}{0.969, 0.969, 0.969}\color{fgcolor}\begin{kframe}
\begin{alltt}
\hldef{B} \hlkwb{<-} \hlnum{1000}
\end{alltt}
\end{kframe}
\end{knitrout}

\begin{knitrout}
\definecolor{shadecolor}{rgb}{0.969, 0.969, 0.969}\color{fgcolor}\begin{kframe}
\begin{alltt}
\hldef{var.monte.carlo} \hlkwb{<-} \hlkwd{c}\hldef{()}

\hlkwa{for} \hldef{(j} \hlkwa{in} \hlnum{1}\hlopt{:}\hlkwd{length}\hldef{(sample.sizes)) \{}

  \hldef{size} \hlkwb{<-} \hldef{sample.sizes[j]}

  \hldef{theta.hat.monte.carlo} \hlkwb{<-} \hlkwd{c}\hldef{()}

  \hlkwa{for} \hldef{(i} \hlkwa{in} \hlnum{1}\hlopt{:}\hldef{B) \{}
    \hldef{u} \hlkwb{<-} \hlkwd{matrix}\hldef{(}\hlkwd{runif}\hldef{(size} \hlopt{*} \hldef{(}\hlnum{2} \hlopt{+} \hlnum{5}\hldef{)),} \hlkwc{nrow} \hldef{= size,} \hlkwc{ncol} \hldef{= (}\hlnum{2} \hlopt{+} \hlnum{5}\hldef{))}

    \hldef{theta.hat.monte.carlo[i]} \hlkwb{<-} \hlkwd{mean}\hldef{(}\hlkwd{my.rbeta}\hldef{(size,} \hlnum{2}\hldef{,} \hlnum{5}\hldef{, u))}
  \hldef{\}}

  \hldef{var.monte.carlo[j]} \hlkwb{<-} \hlkwd{var}\hldef{(theta.hat.monte.carlo)}
\hldef{\}}
\end{alltt}
\end{kframe}
\end{knitrout}

\begin{knitrout}
\definecolor{shadecolor}{rgb}{0.969, 0.969, 0.969}\color{fgcolor}\begin{kframe}
\begin{alltt}
\hldef{var.antithetic} \hlkwb{<-} \hlkwd{c}\hldef{()}

\hlkwa{for} \hldef{(j} \hlkwa{in} \hlnum{1}\hlopt{:}\hlkwd{length}\hldef{(sample.sizes)) \{}

  \hldef{size} \hlkwb{<-} \hldef{sample.sizes[j]}

  \hldef{theta.hat.antithetic} \hlkwb{<-} \hlkwd{c}\hldef{()}

  \hlkwa{for} \hldef{(i} \hlkwa{in} \hlnum{1}\hlopt{:}\hldef{B) \{}
    \hldef{u} \hlkwb{<-} \hlkwd{matrix}\hldef{(}\hlkwd{runif}\hldef{((size} \hlopt{/} \hlnum{2}\hldef{)} \hlopt{*} \hldef{(}\hlnum{2} \hlopt{+} \hlnum{5}\hldef{)),} \hlkwc{nrow} \hldef{= size} \hlopt{/} \hlnum{2}\hldef{,} \hlkwc{ncol} \hldef{= (}\hlnum{2} \hlopt{+} \hlnum{5}\hldef{))}

    \hldef{x} \hlkwb{<-} \hlkwd{my.rbeta}\hldef{(size} \hlopt{/} \hlnum{2}\hldef{,} \hlnum{2}\hldef{,} \hlnum{5}\hldef{, u)}
    \hldef{x_dash} \hlkwb{<-} \hlkwd{my.rbeta}\hldef{(size} \hlopt{/} \hlnum{2}\hldef{,} \hlnum{2}\hldef{,} \hlnum{5}\hldef{,} \hlnum{1}\hlopt{-}\hldef{u)}

    \hldef{theta.hat.antithetic[i]} \hlkwb{<-} \hlkwd{mean}\hldef{(}\hlkwd{c}\hldef{(x, x_dash))}
  \hldef{\}}

  \hldef{var.antithetic[j]} \hlkwb{<-} \hlkwd{var}\hldef{(theta.hat.antithetic)}
\hldef{\}}
\end{alltt}
\end{kframe}
\end{knitrout}

\begin{knitrout}
\definecolor{shadecolor}{rgb}{0.969, 0.969, 0.969}\color{fgcolor}\begin{kframe}
\begin{alltt}
\hldef{relative.gain} \hlkwb{<-} \hldef{((var.monte.carlo} \hlopt{-} \hldef{var.antithetic)} \hlopt{/} \hldef{var.monte.carlo)} \hlopt{*} \hlnum{100}
\end{alltt}
\end{kframe}
\end{knitrout}

\begin{knitrout}
\definecolor{shadecolor}{rgb}{0.969, 0.969, 0.969}\color{fgcolor}\begin{kframe}
\begin{alltt}
\hldef{df2} \hlkwb{<-} \hlkwd{data.frame}\hldef{(}\hlkwc{sample.size} \hldef{= sample.sizes,}
                  \hlkwc{variance.mc} \hldef{= var.monte.carlo,}
                  \hlkwc{variance.at} \hldef{= var.antithetic,}
                  \hlkwc{percentage.reduction} \hldef{= relative.gain)}
\end{alltt}
\end{kframe}
\end{knitrout}

\newpage


% Table created by stargazer v.5.2.3 by Marek Hlavac, Social Policy Institute. E-mail: marek.hlavac at gmail.com
% Date and time: Wed, May 20, 2026 - 02:27:22
\begin{table}[!htbp] \centering 
  \caption{Comparison of Variance of Two Techniques} 
  \label{Table 2} 
\begin{tabular}{@{\extracolsep{5pt}} cccc} 
\\[-1.8ex]\hline 
\hline \\[-1.8ex] 
sample.size & variance.mc & variance.at & percentage.reduction \\ 
\hline \\[-1.8ex] 
$20$ & $0.001220$ & $0.000336$ & $72.428670$ \\ 
$50$ & $0.000507$ & $0.000128$ & $74.805150$ \\ 
$80$ & $0.000331$ & $0.000081$ & $75.702100$ \\ 
$100$ & $0.000252$ & $0.000063$ & $74.983230$ \\ 
\hline \\[-1.8ex] 
\end{tabular} 
\end{table} 


{\fontsize{25}{20} \ding{46}} \hspace{0.1cm} {\fontsize{12}{20}\gloriahallelujahfont We achieve significant reduction in variance of the estimates by using antithetic variable technique.}

\item[\scalebox{2.0}{\twemoji{keycap: 2}}]

\begin{knitrout}
\definecolor{shadecolor}{rgb}{0.969, 0.969, 0.969}\color{fgcolor}\begin{kframe}
\begin{alltt}
\hldef{g} \hlkwb{<-} \hlkwa{function}\hldef{(}\hlkwc{x}\hldef{)} \hlkwd{exp}\hldef{(}\hlopt{-}\hldef{x)}
\end{alltt}
\end{kframe}
\end{knitrout}

\begin{knitrout}
\definecolor{shadecolor}{rgb}{0.969, 0.969, 0.969}\color{fgcolor}\begin{kframe}
\begin{alltt}
\hldef{theta.hat.monte.carlo} \hlkwb{<-} \hlkwd{c}\hldef{()}

\hlkwa{for} \hldef{(i} \hlkwa{in} \hlnum{1}\hlopt{:}\hlkwd{length}\hldef{(sample.sizes)) \{}

  \hldef{size} \hlkwb{<-} \hldef{sample.sizes[i]}

  \hldef{u} \hlkwb{<-} \hlkwd{matrix}\hldef{(}\hlkwd{runif}\hldef{(size} \hlopt{*} \hldef{(}\hlnum{2} \hlopt{+} \hlnum{5}\hldef{)),} \hlkwc{nrow} \hldef{= size,} \hlkwc{ncol} \hldef{= (}\hlnum{2} \hlopt{+} \hlnum{5}\hldef{))}

  \hldef{theta.hat.monte.carlo[i]} \hlkwb{<-} \hlkwd{mean}\hldef{(}\hlkwd{g}\hldef{(}\hlkwd{my.rbeta}\hldef{(size,} \hlnum{2}\hldef{,} \hlnum{5}\hldef{, u)))}
\hldef{\}}
\end{alltt}
\end{kframe}
\end{knitrout}

\begin{knitrout}
\definecolor{shadecolor}{rgb}{0.969, 0.969, 0.969}\color{fgcolor}\begin{kframe}
\begin{alltt}
\hldef{theta.hat.antithetic} \hlkwb{<-} \hlkwd{c}\hldef{()}

\hlkwa{for} \hldef{(i} \hlkwa{in} \hlnum{1}\hlopt{:}\hlkwd{length}\hldef{(sample.sizes)) \{}
  \hldef{size} \hlkwb{<-} \hldef{sample.sizes[i]}

  \hldef{u} \hlkwb{<-} \hlkwd{matrix}\hldef{(}\hlkwd{runif}\hldef{((size} \hlopt{/} \hlnum{2}\hldef{)} \hlopt{*} \hldef{(}\hlnum{2} \hlopt{+} \hlnum{5}\hldef{)),} \hlkwc{nrow} \hldef{= size} \hlopt{/} \hlnum{2}\hldef{,} \hlkwc{ncol} \hldef{= (}\hlnum{2} \hlopt{+} \hlnum{5}\hldef{))}

  \hldef{x} \hlkwb{<-} \hlkwd{my.rbeta}\hldef{(size} \hlopt{/} \hlnum{2}\hldef{,} \hlnum{2}\hldef{,} \hlnum{5}\hldef{, u)}
  \hldef{x_dash} \hlkwb{<-} \hlkwd{my.rbeta}\hldef{(size} \hlopt{/} \hlnum{2}\hldef{,} \hlnum{2}\hldef{,} \hlnum{5}\hldef{,} \hlnum{1}\hlopt{-}\hldef{u)}

  \hldef{theta.hat.antithetic[i]} \hlkwb{<-} \hlkwd{mean}\hldef{(}\hlkwd{g}\hldef{(}\hlkwd{c}\hldef{(x, x_dash)))}
\hldef{\}}
\end{alltt}
\end{kframe}
\end{knitrout}

\begin{knitrout}
\definecolor{shadecolor}{rgb}{0.969, 0.969, 0.969}\color{fgcolor}\begin{kframe}
\begin{alltt}
\hldef{df3} \hlkwb{<-} \hlkwd{data.frame}\hldef{(}\hlkwc{sample.size} \hldef{= sample.sizes,}
                  \hlkwc{theta.hat.mc} \hldef{= theta.hat.monte.carlo,}
                  \hlkwc{theta.hat.at} \hldef{= theta.hat.antithetic)}
\end{alltt}
\end{kframe}
\end{knitrout}


% Table created by stargazer v.5.2.3 by Marek Hlavac, Social Policy Institute. E-mail: marek.hlavac at gmail.com
% Date and time: Wed, May 20, 2026 - 02:27:23
\begin{table}[!htbp] \centering 
  \caption{Estimates from Two Different Techniques} 
  \label{Table 3} 
\begin{tabular}{@{\extracolsep{5pt}} ccc} 
\\[-1.8ex]\hline 
\hline \\[-1.8ex] 
sample.size & theta.hat.mc & theta.hat.at \\ 
\hline \\[-1.8ex] 
$20$ & $0.751372$ & $0.768994$ \\ 
$50$ & $0.775342$ & $0.758296$ \\ 
$80$ & $0.758218$ & $0.758144$ \\ 
$100$ & $0.752671$ & $0.752472$ \\ 
\hline \\[-1.8ex] 
\end{tabular} 
\end{table} 


\begin{knitrout}
\definecolor{shadecolor}{rgb}{0.969, 0.969, 0.969}\color{fgcolor}\begin{kframe}
\begin{alltt}
\hldef{B} \hlkwb{<-} \hlnum{1000}
\end{alltt}
\end{kframe}
\end{knitrout}

\begin{knitrout}
\definecolor{shadecolor}{rgb}{0.969, 0.969, 0.969}\color{fgcolor}\begin{kframe}
\begin{alltt}
\hldef{var.monte.carlo} \hlkwb{<-} \hlkwd{c}\hldef{()}

\hlkwa{for} \hldef{(j} \hlkwa{in} \hlnum{1}\hlopt{:}\hlkwd{length}\hldef{(sample.sizes)) \{}

  \hldef{size} \hlkwb{<-} \hldef{sample.sizes[j]}

  \hldef{theta.hat.monte.carlo} \hlkwb{<-} \hlkwd{c}\hldef{()}

  \hlkwa{for} \hldef{(i} \hlkwa{in} \hlnum{1}\hlopt{:}\hldef{B) \{}
    \hldef{u} \hlkwb{<-} \hlkwd{matrix}\hldef{(}\hlkwd{runif}\hldef{(size} \hlopt{*} \hldef{(}\hlnum{2} \hlopt{+} \hlnum{5}\hldef{)),} \hlkwc{nrow} \hldef{= size,} \hlkwc{ncol} \hldef{= (}\hlnum{2} \hlopt{+} \hlnum{5}\hldef{))}

    \hldef{theta.hat.monte.carlo[i]} \hlkwb{<-} \hlkwd{mean}\hldef{(}\hlkwd{g}\hldef{(}\hlkwd{my.rbeta}\hldef{(size,} \hlnum{2}\hldef{,} \hlnum{5}\hldef{, u)))}
  \hldef{\}}

  \hldef{var.monte.carlo[j]} \hlkwb{<-} \hlkwd{var}\hldef{(theta.hat.monte.carlo)}
\hldef{\}}
\end{alltt}
\end{kframe}
\end{knitrout}

\begin{knitrout}
\definecolor{shadecolor}{rgb}{0.969, 0.969, 0.969}\color{fgcolor}\begin{kframe}
\begin{alltt}
\hldef{var.antithetic} \hlkwb{<-} \hlkwd{c}\hldef{()}

\hlkwa{for} \hldef{(j} \hlkwa{in} \hlnum{1}\hlopt{:}\hlkwd{length}\hldef{(sample.sizes)) \{}

  \hldef{size} \hlkwb{<-} \hldef{sample.sizes[j]}

  \hldef{theta.hat.antithetic} \hlkwb{<-} \hlkwd{c}\hldef{()}

  \hlkwa{for} \hldef{(i} \hlkwa{in} \hlnum{1}\hlopt{:}\hldef{B) \{}
    \hldef{u} \hlkwb{<-} \hlkwd{matrix}\hldef{(}\hlkwd{runif}\hldef{((size} \hlopt{/} \hlnum{2}\hldef{)} \hlopt{*} \hldef{(}\hlnum{2} \hlopt{+} \hlnum{5}\hldef{)),} \hlkwc{nrow} \hldef{= size} \hlopt{/} \hlnum{2}\hldef{,} \hlkwc{ncol} \hldef{= (}\hlnum{2} \hlopt{+} \hlnum{5}\hldef{))}

    \hldef{x} \hlkwb{<-} \hlkwd{my.rbeta}\hldef{(size} \hlopt{/} \hlnum{2}\hldef{,} \hlnum{2}\hldef{,} \hlnum{5}\hldef{, u)}
    \hldef{x_dash} \hlkwb{<-} \hlkwd{my.rbeta}\hldef{(size} \hlopt{/} \hlnum{2}\hldef{,} \hlnum{2}\hldef{,} \hlnum{5}\hldef{,} \hlnum{1}\hlopt{-}\hldef{u)}

    \hldef{theta.hat.antithetic[i]} \hlkwb{<-} \hlkwd{mean}\hldef{(}\hlkwd{g}\hldef{(}\hlkwd{c}\hldef{(x, x_dash)))}
  \hldef{\}}

  \hldef{var.antithetic[j]} \hlkwb{<-} \hlkwd{var}\hldef{(theta.hat.antithetic)}
\hldef{\}}
\end{alltt}
\end{kframe}
\end{knitrout}

\begin{knitrout}
\definecolor{shadecolor}{rgb}{0.969, 0.969, 0.969}\color{fgcolor}\begin{kframe}
\begin{alltt}
\hldef{relative.gain} \hlkwb{<-} \hldef{((var.monte.carlo} \hlopt{-} \hldef{var.antithetic)} \hlopt{/} \hldef{var.monte.carlo)} \hlopt{*} \hlnum{100}
\end{alltt}
\end{kframe}
\end{knitrout}

\begin{knitrout}
\definecolor{shadecolor}{rgb}{0.969, 0.969, 0.969}\color{fgcolor}\begin{kframe}
\begin{alltt}
\hldef{df4} \hlkwb{<-} \hlkwd{data.frame}\hldef{(}\hlkwc{sample.size} \hldef{= sample.sizes,}
                  \hlkwc{variance.mc} \hldef{= var.monte.carlo,}
                  \hlkwc{variance.at} \hldef{= var.antithetic,}
                  \hlkwc{percentage.reduction} \hldef{= relative.gain)}
\end{alltt}
\end{kframe}
\end{knitrout}

\newpage


% Table created by stargazer v.5.2.3 by Marek Hlavac, Social Policy Institute. E-mail: marek.hlavac at gmail.com
% Date and time: Wed, May 20, 2026 - 02:27:29
\begin{table}[!htbp] \centering 
  \caption{Comparison of Variance of Two Techniques} 
  \label{Table 4} 
\begin{tabular}{@{\extracolsep{5pt}} cccc} 
\\[-1.8ex]\hline 
\hline \\[-1.8ex] 
sample.size & variance.mc & variance.at & percentage.reduction \\ 
\hline \\[-1.8ex] 
$20$ & $0.000703$ & $0.000132$ & $81.251030$ \\ 
$50$ & $0.000257$ & $0.000057$ & $77.672460$ \\ 
$80$ & $0.000159$ & $0.000037$ & $76.891230$ \\ 
$100$ & $0.000127$ & $0.000031$ & $75.728570$ \\ 
\hline \\[-1.8ex] 
\end{tabular} 
\end{table} 


{\fontsize{25}{20} \ding{46}} \hspace{0.1cm} {\fontsize{12}{20}\gloriahallelujahfont We achieve significant reduction in variance of the estimates by using antithetic variable technique.}



\item[\scalebox{2.0}{\twemoji{keycap: 3}}]

\begin{knitrout}
\definecolor{shadecolor}{rgb}{0.969, 0.969, 0.969}\color{fgcolor}\begin{kframe}
\begin{alltt}
\hldef{g} \hlkwb{<-} \hlkwa{function}\hldef{(}\hlkwc{x}\hldef{) x}\hlopt{^}\hlnum{2}
\end{alltt}
\end{kframe}
\end{knitrout}

\begin{knitrout}
\definecolor{shadecolor}{rgb}{0.969, 0.969, 0.969}\color{fgcolor}\begin{kframe}
\begin{alltt}
\hldef{theta.hat.monte.carlo} \hlkwb{<-} \hlkwd{c}\hldef{()}

\hlkwa{for} \hldef{(i} \hlkwa{in} \hlnum{1}\hlopt{:}\hlkwd{length}\hldef{(sample.sizes)) \{}

  \hldef{size} \hlkwb{<-} \hldef{sample.sizes[i]}

  \hldef{u} \hlkwb{<-} \hlkwd{matrix}\hldef{(}\hlkwd{runif}\hldef{(size} \hlopt{*} \hldef{(}\hlnum{2} \hlopt{+} \hlnum{5}\hldef{)),} \hlkwc{nrow} \hldef{= size,} \hlkwc{ncol} \hldef{= (}\hlnum{2} \hlopt{+} \hlnum{5}\hldef{))}

  \hldef{theta.hat.monte.carlo[i]} \hlkwb{<-} \hlkwd{mean}\hldef{(}\hlkwd{g}\hldef{(}\hlkwd{my.rbeta}\hldef{(size,} \hlnum{2}\hldef{,} \hlnum{5}\hldef{, u)))}
\hldef{\}}
\end{alltt}
\end{kframe}
\end{knitrout}

\begin{knitrout}
\definecolor{shadecolor}{rgb}{0.969, 0.969, 0.969}\color{fgcolor}\begin{kframe}
\begin{alltt}
\hldef{theta.hat.antithetic} \hlkwb{<-} \hlkwd{c}\hldef{()}

\hlkwa{for} \hldef{(i} \hlkwa{in} \hlnum{1}\hlopt{:}\hlkwd{length}\hldef{(sample.sizes)) \{}
  \hldef{size} \hlkwb{<-} \hldef{sample.sizes[i]}

  \hldef{u} \hlkwb{<-} \hlkwd{matrix}\hldef{(}\hlkwd{runif}\hldef{((size} \hlopt{/} \hlnum{2}\hldef{)} \hlopt{*} \hldef{(}\hlnum{2} \hlopt{+} \hlnum{5}\hldef{)),} \hlkwc{nrow} \hldef{= size} \hlopt{/} \hlnum{2}\hldef{,} \hlkwc{ncol} \hldef{= (}\hlnum{2} \hlopt{+} \hlnum{5}\hldef{))}

  \hldef{x} \hlkwb{<-} \hlkwd{my.rbeta}\hldef{(size} \hlopt{/} \hlnum{2}\hldef{,} \hlnum{2}\hldef{,} \hlnum{5}\hldef{, u)}
  \hldef{x_dash} \hlkwb{<-} \hlkwd{my.rbeta}\hldef{(size} \hlopt{/} \hlnum{2}\hldef{,} \hlnum{2}\hldef{,} \hlnum{5}\hldef{,} \hlnum{1}\hlopt{-}\hldef{u)}

  \hldef{theta.hat.antithetic[i]} \hlkwb{<-} \hlkwd{mean}\hldef{(}\hlkwd{g}\hldef{(}\hlkwd{c}\hldef{(x, x_dash)))}
\hldef{\}}
\end{alltt}
\end{kframe}
\end{knitrout}

\begin{knitrout}
\definecolor{shadecolor}{rgb}{0.969, 0.969, 0.969}\color{fgcolor}\begin{kframe}
\begin{alltt}
\hldef{df5} \hlkwb{<-} \hlkwd{data.frame}\hldef{(}\hlkwc{sample.size} \hldef{= sample.sizes,}
                  \hlkwc{theta.hat.mc} \hldef{= theta.hat.monte.carlo,}
                  \hlkwc{theta.hat.at} \hldef{= theta.hat.antithetic)}
\end{alltt}
\end{kframe}
\end{knitrout}


% Table created by stargazer v.5.2.3 by Marek Hlavac, Social Policy Institute. E-mail: marek.hlavac at gmail.com
% Date and time: Wed, May 20, 2026 - 02:27:29
\begin{table}[!htbp] \centering 
  \caption{Estimates from Two Different Techniques} 
  \label{Table 5} 
\begin{tabular}{@{\extracolsep{5pt}} ccc} 
\\[-1.8ex]\hline 
\hline \\[-1.8ex] 
sample.size & theta.hat.mc & theta.hat.at \\ 
\hline \\[-1.8ex] 
$20$ & $0.126081$ & $0.096772$ \\ 
$50$ & $0.118958$ & $0.105170$ \\ 
$80$ & $0.117715$ & $0.112296$ \\ 
$100$ & $0.096335$ & $0.109589$ \\ 
\hline \\[-1.8ex] 
\end{tabular} 
\end{table} 


\newpage

True value is

\begin{align*}
\theta &= E[X^2] \\[0.2em]
&= \displaystyle{\int \limits_0^1 x^2 \cdot \dfrac{1}{\mathcal{B}(2, 5)} x (1 - x)^4 \; dx} \\[0.5em]
&= \displaystyle{\dfrac{1}{\mathcal{B}(2, 5)} \int \limits_0^1 x^3 (1 - x)^4 \; dx} \\[0.5em]
&= \dfrac{1}{\mathcal{B}(2, 5)} \cdot \mathcal{B}(4, 5) \\[0.5em]
&= \dfrac{\Gamma(7)}{\Gamma(2) \cdot \Gamma(5)} \cdot \dfrac{\Gamma(4) \cdot \Gamma(5)}{\Gamma(9)} \\[0.5em]
&= \dfrac{6!}{4!} \cdot \dfrac{3! \cdot 4!}{8!} \\[0.5em]
&= \dfrac{6}{7 \times 8} \\[0.5em]
&= \dfrac{3}{28} \\[0.5em]
&= 0.1071429.
\end{align*}

\begin{knitrout}
\definecolor{shadecolor}{rgb}{0.969, 0.969, 0.969}\color{fgcolor}\begin{kframe}
\begin{alltt}
\hldef{B} \hlkwb{<-} \hlnum{1000}
\end{alltt}
\end{kframe}
\end{knitrout}

\begin{knitrout}
\definecolor{shadecolor}{rgb}{0.969, 0.969, 0.969}\color{fgcolor}\begin{kframe}
\begin{alltt}
\hldef{var.monte.carlo} \hlkwb{<-} \hlkwd{c}\hldef{()}

\hlkwa{for} \hldef{(j} \hlkwa{in} \hlnum{1}\hlopt{:}\hlkwd{length}\hldef{(sample.sizes)) \{}

  \hldef{size} \hlkwb{<-} \hldef{sample.sizes[j]}

  \hldef{theta.hat.monte.carlo} \hlkwb{<-} \hlkwd{c}\hldef{()}

  \hlkwa{for} \hldef{(i} \hlkwa{in} \hlnum{1}\hlopt{:}\hldef{B) \{}
    \hldef{u} \hlkwb{<-} \hlkwd{matrix}\hldef{(}\hlkwd{runif}\hldef{(size} \hlopt{*} \hldef{(}\hlnum{2} \hlopt{+} \hlnum{5}\hldef{)),} \hlkwc{nrow} \hldef{= size,} \hlkwc{ncol} \hldef{= (}\hlnum{2} \hlopt{+} \hlnum{5}\hldef{))}

    \hldef{theta.hat.monte.carlo[i]} \hlkwb{<-} \hlkwd{mean}\hldef{(}\hlkwd{g}\hldef{(}\hlkwd{my.rbeta}\hldef{(size,} \hlnum{2}\hldef{,} \hlnum{5}\hldef{, u)))}
  \hldef{\}}

  \hldef{var.monte.carlo[j]} \hlkwb{<-} \hlkwd{var}\hldef{(theta.hat.monte.carlo)}
\hldef{\}}
\end{alltt}
\end{kframe}
\end{knitrout}

\begin{knitrout}
\definecolor{shadecolor}{rgb}{0.969, 0.969, 0.969}\color{fgcolor}\begin{kframe}
\begin{alltt}
\hldef{var.antithetic} \hlkwb{<-} \hlkwd{c}\hldef{()}

\hlkwa{for} \hldef{(j} \hlkwa{in} \hlnum{1}\hlopt{:}\hlkwd{length}\hldef{(sample.sizes)) \{}

  \hldef{size} \hlkwb{<-} \hldef{sample.sizes[j]}

  \hldef{theta.hat.antithetic} \hlkwb{<-} \hlkwd{c}\hldef{()}

  \hlkwa{for} \hldef{(i} \hlkwa{in} \hlnum{1}\hlopt{:}\hldef{B) \{}
    \hldef{u} \hlkwb{<-} \hlkwd{matrix}\hldef{(}\hlkwd{runif}\hldef{((size} \hlopt{/} \hlnum{2}\hldef{)} \hlopt{*} \hldef{(}\hlnum{2} \hlopt{+} \hlnum{5}\hldef{)),} \hlkwc{nrow} \hldef{= size} \hlopt{/} \hlnum{2}\hldef{,} \hlkwc{ncol} \hldef{= (}\hlnum{2} \hlopt{+} \hlnum{5}\hldef{))}

    \hldef{x} \hlkwb{<-} \hlkwd{my.rbeta}\hldef{(size} \hlopt{/} \hlnum{2}\hldef{,} \hlnum{2}\hldef{,} \hlnum{5}\hldef{, u)}
    \hldef{x_dash} \hlkwb{<-} \hlkwd{my.rbeta}\hldef{(size} \hlopt{/} \hlnum{2}\hldef{,} \hlnum{2}\hldef{,} \hlnum{5}\hldef{,} \hlnum{1}\hlopt{-}\hldef{u)}

    \hldef{theta.hat.antithetic[i]} \hlkwb{<-} \hlkwd{mean}\hldef{(}\hlkwd{g}\hldef{(}\hlkwd{c}\hldef{(x, x_dash)))}
  \hldef{\}}

  \hldef{var.antithetic[j]} \hlkwb{<-} \hlkwd{var}\hldef{(theta.hat.antithetic)}
\hldef{\}}
\end{alltt}
\end{kframe}
\end{knitrout}

\begin{knitrout}
\definecolor{shadecolor}{rgb}{0.969, 0.969, 0.969}\color{fgcolor}\begin{kframe}
\begin{alltt}
\hldef{relative.gain} \hlkwb{<-} \hldef{((var.monte.carlo} \hlopt{-} \hldef{var.antithetic)} \hlopt{/} \hldef{var.monte.carlo)} \hlopt{*} \hlnum{100}
\end{alltt}
\end{kframe}
\end{knitrout}

\begin{knitrout}
\definecolor{shadecolor}{rgb}{0.969, 0.969, 0.969}\color{fgcolor}\begin{kframe}
\begin{alltt}
\hldef{df6} \hlkwb{<-} \hlkwd{data.frame}\hldef{(}\hlkwc{sample.size} \hldef{= sample.sizes,}
                  \hlkwc{variance.mc} \hldef{= var.monte.carlo,}
                  \hlkwc{variance.at} \hldef{= var.antithetic,}
                  \hlkwc{percentage.reduction} \hldef{= relative.gain)}
\end{alltt}
\end{kframe}
\end{knitrout}


% Table created by stargazer v.5.2.3 by Marek Hlavac, Social Policy Institute. E-mail: marek.hlavac at gmail.com
% Date and time: Wed, May 20, 2026 - 02:27:33
\begin{table}[!htbp] \centering 
  \caption{Comparison of Variance of Two Techniques} 
  \label{Table 6} 
\begin{tabular}{@{\extracolsep{5pt}} cccc} 
\\[-1.8ex]\hline 
\hline \\[-1.8ex] 
sample.size & variance.mc & variance.at & percentage.reduction \\ 
\hline \\[-1.8ex] 
$20$ & $0.000617$ & $0.000287$ & $53.495410$ \\ 
$50$ & $0.000271$ & $0.000127$ & $53.163210$ \\ 
$80$ & $0.000153$ & $0.000070$ & $54.183050$ \\ 
$100$ & $0.000125$ & $0.000054$ & $56.463350$ \\ 
\hline \\[-1.8ex] 
\end{tabular} 
\end{table} 


{\fontsize{25}{20} \ding{46}} \hspace{0.1cm} {\fontsize{12}{20}\gloriahallelujahfont We achieve significant reduction in variance of the estimates by using antithetic variable technique.}

\end{enumerate}

\end{document}
